% ============================================================================
% RPS.TEX
% Rencana Pembelajaran Semester (RPS) Terintegrasi
% Dengan pemetaan CPL-CPMK-Bab-Section
% ============================================================================

\chapter{Rencana Pembelajaran Semester (RPS)}

\section{Identitas Mata Kuliah}

\begin{table}[h]
\centering
\begin{tabular}{ll}
\textbf{Nama Mata Kuliah} & : Pemrograman Game \\
\textbf{Kode Mata Kuliah} & : IF-XXX \\
\textbf{Bobot SKS} & : 3 SKS \\
\textbf{Semester} & : [Tentukan Semester] \\
\textbf{Program Studi} & : S1 Informatika \\
\textbf{Dosen Pengampu} & : [Nama Dosen] \\
\textbf{Prasyarat} & : Pemrograman Berorientasi Objek \\
\end{tabular}
\end{table}

\section{Deskripsi Mata Kuliah}

Mata kuliah Pemrograman Game merupakan mata kuliah yang membekali mahasiswa dengan pengetahuan dan keterampilan praktis dalam pengembangan game menggunakan Unity Engine dan bahasa pemrograman C\#. Mata kuliah ini mengadopsi pendekatan Outcome-Based Education (OBE) dengan fokus pada pembelajaran berbasis proyek (project-based learning) yang dominan praktik dibanding teori.

Mata kuliah ini mencakup konsep dasar game development dan game design, pengenalan Unity Engine dan workflow pengembangan game, pemrograman C\# untuk Unity, pengelolaan game object dan komponen, sistem input dan kontrol pemain, fisika dan collision detection, animasi, UI/UX game, audio, level design, AI sederhana, optimasi, serta proses build dan deployment game. Melalui pembelajaran bertahap dan proyek akhir yang komprehensif, mahasiswa diharapkan mampu merancang, mengembangkan, menguji, dan mempublikasikan game yang fungsional dan menarik.

\section{Capaian Pembelajaran Mata Kuliah (CPMK)}

Setelah menyelesaikan mata kuliah ini, mahasiswa diharapkan mampu:

\begin{enumerate}[label=CPMK\arabic*:]
\item Menganalisis konsep dasar game development, game design, dan genre game serta menjelaskan workflow pengembangan game menggunakan Unity Engine.

\item Menerapkan pemrograman C\# untuk Unity dengan menguasai struktur skrip, fungsi dasar (\texttt{Start()}, \texttt{Update()}), variabel, kontrol alur program, dan konsep object-oriented programming dalam konteks game development.

\item Mengoperasikan Unity Editor dengan mahir dalam mengelola game object, komponen (Transform, Rigidbody, Collider), prefab, scene, dan manajemen aset (grafik, model 3D, audio).

\item Mengimplementasikan sistem kontrol pemain menggunakan Input System Unity, menggerakkan karakter, dan menangani interaksi pemain dengan game world.

\item Mengaplikasikan sistem fisika Unity (Rigidbody, Collider, Physics Material) untuk simulasi gerak, gravitasi, tumbukan, dan interaksi fisik dalam game.

\item Membuat animasi karakter dan objek menggunakan Animator Controller, mengelola transisi animasi, dan mengintegrasikan animasi dengan gameplay.

\item Merancang dan mengimplementasikan UI/UX game yang efektif meliputi HUD (Heads-Up Display), menu navigasi, sistem scoring, dan feedback visual untuk pengalaman pengguna yang optimal.

\item Mengintegrasikan audio dalam game meliputi sound effect dan background music menggunakan AudioSource, AudioListener, dan Audio Mixer untuk meningkatkan immersi gameplay.

\item Merancang level design dan game mechanics yang menarik, termasuk sistem progresi, tantangan, dan reward yang seimbang.

\item Mengimplementasikan AI sederhana untuk NPC (Non-Player Character) menggunakan teknik dasar seperti pathfinding dengan NavMesh, state machine, dan behavior scripting.

\item Menganalisis dan mengoptimalkan performa game melalui profiling, debugging, dan teknik optimasi aset serta kode untuk berbagai platform target.

\item Membangun dan mendeploy game ke berbagai platform (PC, Mobile, Web) dengan memahami proses build, pengaturan platform-specific, dan distribusi game.

\item Merancang dan mengembangkan proyek game utuh yang mencakup semua aspek pembelajaran mulai dari konsep, desain, implementasi, pengujian, hingga publikasi dengan menerapkan best practices industri game development.
\end{enumerate}

\section{Pemetaan CPL-CPMK-Bab-Section}

\begin{longtable}{|p{1cm}|p{2cm}|p{3cm}|p{2.5cm}|p{3cm}|}
\hline
\textbf{CPL} & \textbf{CPMK} & \textbf{Bab} & \textbf{Section} & \textbf{Indikator Kinerja} \\
\hline
\endfirsthead

\hline
\textbf{CPL} & \textbf{CPMK} & \textbf{Bab} & \textbf{Section} & \textbf{Indikator Kinerja} \\
\hline
\endhead

\hline
\endfoot

\hline
\endlastfoot

CPL1, CPL2 & CPMK1 & Bab 1 & 1.1-1.5 & Mampu menjelaskan konsep game dan workflow Unity \\
CPL3, CPL4 & CPMK2 & Bab 3 & 3.1-3.7 & Mampu membuat skrip C\# untuk Unity \\
CPL4 & CPMK3 & Bab 2, 4 & 2.1-2.5, 4.1-4.5 & Mampu mengoperasikan Unity Editor \\
CPL3, CPL4 & CPMK4 & Bab 5 & 5.1-5.6 & Mampu implementasi input system \\
CPL2, CPL4 & CPMK5 & Bab 6 & 6.1-6.6 & Mampu aplikasi sistem fisika \\
CPL4 & CPMK6 & Bab 7 & 7.1-7.5 & Mampu membuat animasi \\
CPL3, CPL4 & CPMK7 & Bab 8 & 8.1-8.5 & Mampu desain UI/UX \\
CPL4 & CPMK8 & Bab 9 & 9.1-9.4 & Mampu integrasi audio \\
CPL2, CPL3 & CPMK9 & Bab 10 & 10.1-10.5 & Mampu desain level \\
CPL2, CPL4 & CPMK10 & Bab 11 & 11.1-11.5 & Mampu implementasi AI \\
CPL3, CPL4 & CPMK11 & Bab 12 & 12.1-12.5 & Mampu optimasi performa \\
CPL3, CPL4 & CPMK12 & Bab 13 & 13.1-13.5 & Mampu build dan deploy \\
CPL1, CPL3, CPL4 & CPMK13 & Semua Bab & Semua Section & Mampu membuat game utuh \\
\hline
\end{longtable}

\section{Rencana Pembelajaran Semester (RPS) 16 Pertemuan}

\begin{longtable}{|p{0.8cm}|p{3.5cm}|p{4cm}|p{4cm}|p{2.5cm}|}
\hline
\textbf{No} & \textbf{Tujuan Pembelajaran} & \textbf{Materi Pembelajaran} & \textbf{Aktivitas Praktik} & \textbf{Evaluasi} \\
\hline
\endfirsthead

\hline
\textbf{No} & \textbf{Tujuan Pembelajaran} & \textbf{Materi Pembelajaran} & \textbf{Aktivitas Praktik} & \textbf{Evaluasi} \\
\hline
\endhead

\hline
\endfoot

\hline
\endlastfoot

1 & Mahasiswa mampu menjelaskan konsep dasar game development, genre game, dan workflow pengembangan game menggunakan Unity. & 
\begin{itemize}[leftmargin=*,nosep]
\item Konsep dasar game dan elemen-elemen game
\item Genre game dan karakteristiknya
\item Lifecycle pengembangan game
\item Pengenalan Unity Engine dan keunggulannya
\item Instalasi Unity Hub dan Unity Editor
\end{itemize} & 
\begin{itemize}[leftmargin=*,nosep]
\item Diskusi tentang game favorit dan analisis elemen-elemennya
\item Instalasi Unity Hub dan Unity Editor
\item Eksplorasi Unity Editor interface
\item Membuat project Unity pertama
\end{itemize} & 
Kuis: Konsep dasar game development \\
\hline

2 & Mahasiswa mampu mengoperasikan Unity Editor dengan mahir dan memahami struktur project Unity. & 
\begin{itemize}[leftmargin=*,nosep]
\item Unity Editor interface detail (Scene, Game, Hierarchy, Inspector, Project, Console)
\item Workflow pengembangan game di Unity
\item Membuat dan mengelola Scene
\item Unity Asset Store dan package management
\end{itemize} & 
\begin{itemize}[leftmargin=*,nosep]
\item Praktik navigasi Unity Editor
\item Membuat Scene baru dengan berbagai objek dasar
\item Import asset dari Asset Store
\item Setup project structure yang rapi
\end{itemize} & 
Tugas: Setup project Unity dengan struktur folder yang baik \\
\hline

3 & Mahasiswa mampu membuat skrip C\# dasar untuk Unity dan memahami lifecycle MonoBehaviour. & 
\begin{itemize}[leftmargin=*,nosep]
\item Struktur skrip C\# di Unity (MonoBehaviour)
\item Fungsi \texttt{Start()}, \texttt{Update()}, \texttt{Awake()}, \texttt{OnEnable()}
\item Variabel dan tipe data dasar
\item Debug.Log untuk debugging
\end{itemize} & 
\begin{itemize}[leftmargin=*,nosep]
\item Membuat skrip C\# pertama di Unity
\item Praktik penggunaan fungsi lifecycle
\item Manipulasi variabel dan output ke Console
\item Latihan membuat skrip sederhana (rotasi objek)
\end{itemize} & 
Tugas: Membuat skrip rotasi dan pergerakan objek sederhana \\
\hline

4 & Mahasiswa mampu menguasai konsep GameObject, Component, Transform, dan Prefab. & 
\begin{itemize}[leftmargin=*,nosep]
\item Konsep GameObject dan Hierarchy
\item Transform component (position, rotation, scale)
\item Komponen-komponen Unity dasar
\item Prefab dan instantiation
\item Tag dan Layer
\end{itemize} & 
\begin{itemize}[leftmargin=*,nosep]
\item Praktik manipulasi Transform melalui Inspector dan script
\item Membuat Prefab dari GameObject
\item Instantiate Prefab melalui script
\item Latihan membuat scene dengan berbagai objek dan prefab
\end{itemize} & 
Tugas: Membuat scene dengan prefab system \\
\hline

5 & Mahasiswa mampu mengimplementasikan sistem input dan kontrol pemain dasar. & 
\begin{itemize}[leftmargin=*,nosep]
\item Unity Input System (legacy dan new)
\item Keyboard dan mouse input
\item Player movement 2D dan 3D
\item Camera control dasar
\end{itemize} & 
\begin{itemize}[leftmargin=*,nosep]
\item Implementasi player movement dengan keyboard
\item Membuat karakter yang dapat bergerak (WASD/Arrow keys)
\item Setup camera follow untuk karakter
\item Latihan membuat kontrol pemain yang responsif
\end{itemize} & 
Tugas: Game sederhana dengan player controller (contoh: menggerakkan kotak) \\
\hline

6 & Mahasiswa mampu mengaplikasikan sistem fisika Unity untuk simulasi gerak dan tumbukan. & 
\begin{itemize}[leftmargin=*,nosep]
\item Rigidbody component dan propertinya
\item Collider types dan collision detection
\item Physics Material
\item OnCollisionEnter dan OnTriggerEnter
\end{itemize} & 
\begin{itemize}[leftmargin=*,nosep]
\item Setup Rigidbody dan Collider pada objek
\item Membuat simulasi fisika (jatuh, tumbukan)
\item Implementasi trigger untuk collectibles
\item Latihan membuat game dengan mekanik fisika (contoh: platformer sederhana)
\end{itemize} & 
Tugas: Game dengan sistem fisika (contoh: platformer atau puzzle game sederhana) \\
\hline

7 & Mahasiswa mampu membuat animasi karakter menggunakan Animator Controller. & 
\begin{itemize}[leftmargin=*,nosep]
\item Konsep animasi dalam game
\item Animator Controller dan state machine
\item Animation states dan transitions
\item Parameters (float, int, bool, trigger)
\end{itemize} & 
\begin{itemize}[leftmargin=*,nosep]
\item Membuat Animator Controller
\item Setup animation states (idle, walk, jump)
\item Membuat transitions antar states
\item Kontrol animasi melalui script
\end{itemize} & 
Tugas: Karakter dengan animasi dasar (idle, walk, jump) \\
\hline

8 & \textbf{UTS - Ujian Tengah Semester} & 
\begin{itemize}[leftmargin=*,nosep]
\item Review materi pertemuan 1-7
\item Konsep game development
\item Unity Editor dan workflow
\item Pemrograman C\# untuk Unity
\item GameObject, Component, Prefab
\item Input System dan Player Controller
\item Physics dan Collision
\item Animation dasar
\end{itemize} & 
\begin{itemize}[leftmargin=*,nosep]
\item Ujian teori (pilihan ganda dan esai)
\item Presentasi progress proyek game individu
\item Demo game sederhana yang telah dibuat
\end{itemize} & 
UTS: Ujian teori dan presentasi proyek \\
\hline

9 & Mahasiswa mampu merancang dan mengimplementasikan UI/UX game yang efektif. & 
\begin{itemize}[leftmargin=*,nosep]
\item Canvas dan Canvas Scaler
\item UI Components (Button, Text, Image, Slider)
\item Event System dan UI events
\item HUD design (health bar, score display)
\item Menu system (Main menu, Pause menu)
\end{itemize} & 
\begin{itemize}[leftmargin=*,nosep]
\item Membuat Canvas dan setup UI
\item Implementasi HUD (score, health bar)
\item Membuat Main Menu dengan navigasi
\item Implementasi Pause Menu
\item Latihan membuat UI yang responsif
\end{itemize} & 
Tugas: Game dengan UI lengkap (HUD + Menu) \\
\hline

10 & Mahasiswa mampu mengintegrasikan audio dalam game untuk meningkatkan immersi. & 
\begin{itemize}[leftmargin=*,nosep]
\item AudioSource dan AudioListener
\item Sound effect vs background music
\item 2D vs 3D audio
\item Audio Mixer dasar
\item Dynamic audio control
\end{itemize} & 
\begin{itemize}[leftmargin=*,nosep]
\item Import audio assets
\item Setup AudioSource untuk sound effect
\item Implementasi background music
\item Kontrol volume melalui UI
\item Latihan membuat audio yang sesuai dengan gameplay
\end{itemize} & 
Tugas: Menambahkan audio (SFX + BGM) ke game \\
\hline

11 & Mahasiswa mampu merancang level design dan game mechanics yang menarik. & 
\begin{itemize}[leftmargin=*,nosep]
\item Prinsip level design (flow, pacing, challenge)
\item Game state management
\item Win/lose conditions
\item Spawn system dan object pooling
\item Collectibles dan power-ups
\end{itemize} & 
\begin{itemize}[leftmargin=*,nosep]
\item Merancang level dengan prinsip level design
\item Implementasi game state (playing, paused, game over)
\item Membuat sistem spawn untuk enemies/objects
\item Implementasi collectibles dan scoring
\item Latihan membuat level yang seimbang dan menarik
\end{itemize} & 
Tugas: Level design dengan mechanics lengkap \\
\hline

12 & Mahasiswa mampu mengimplementasikan AI sederhana untuk NPC menggunakan teknik dasar. & 
\begin{itemize}[leftmargin=*,nosep]
\item Konsep AI dalam game
\item NPC behavior patterns (patrol, chase, attack)
\item NavMesh dan pathfinding
\item State Machine untuk AI
\item Line of sight detection
\end{itemize} & 
\begin{itemize}[leftmargin=*,nosep]
\item Setup NavMesh untuk pathfinding
\item Membuat NPC dengan behavior patrol
\item Implementasi chase behavior untuk enemy
\item State Machine untuk AI (idle, patrol, chase, attack)
\item Latihan membuat AI yang menantang namun fair
\end{itemize} & 
Tugas: Game dengan AI enemy (patrol + chase) \\
\hline

13 & Mahasiswa mampu menganalisis dan mengoptimalkan performa game. & 
\begin{itemize}[leftmargin=*,nosep]
\item Unity Profiler dan Frame Debugger
\item Performance bottlenecks (CPU, GPU, memory)
\item Optimization techniques (object pooling, LOD, occlusion culling)
\item Debugging tools dan techniques
\end{itemize} & 
\begin{itemize}[leftmargin=*,nosep]
\item Menggunakan Profiler untuk analisis performa
\item Implementasi object pooling
\item Optimasi texture dan model
\item Debugging game yang sudah dibuat
\item Latihan mengoptimalkan game untuk performa lebih baik
\end{itemize} & 
Tugas: Optimasi game yang telah dibuat \\
\hline

14 & Mahasiswa mampu membangun dan mendeploy game ke berbagai platform. & 
\begin{itemize}[leftmargin=*,nosep]
\item Build settings dan platform selection
\item Platform-specific settings (PC, Android, WebGL)
\item Asset optimization untuk target platform
\item Build process dan troubleshooting
\item Distribution channels
\end{itemize} & 
\begin{itemize}[leftmargin=*,nosep]
\item Setup build settings untuk PC
\item Build game untuk WebGL
\item Setup build untuk Android (jika memungkinkan)
\item Testing build pada target platform
\item Upload ke platform distribusi (itch.io)
\end{itemize} & 
Tugas: Build dan deploy game ke platform \\
\hline

15 & Mahasiswa mampu mengembangkan proyek game utuh dengan menerapkan semua aspek pembelajaran. & 
\begin{itemize}[leftmargin=*,nosep]
\item Review semua materi pembelajaran
\item Best practices game development
\item Project management untuk game development
\item Version control dengan Git
\item Testing dan quality assurance
\end{itemize} & 
\begin{itemize}[leftmargin=*,nosep]
\item Finalisasi proyek game akhir
\item Implementasi fitur-fitur yang belum ada
\item Testing menyeluruh
\item Dokumentasi proyek
\item Persiapan presentasi final
\end{itemize} & 
Progress check: Review proyek akhir \\
\hline

16 & \textbf{UAS - Ujian Akhir Semester} & 
\begin{itemize}[leftmargin=*,nosep]
\item Presentasi proyek game akhir
\item Demo game yang telah dikembangkan
\item Q\&A tentang implementasi dan konsep
\item Refleksi pembelajaran
\end{itemize} & 
\begin{itemize}[leftmargin=*,nosep]
\item Presentasi proyek game akhir (15-20 menit)
\item Demo gameplay
\item Penjelasan implementasi teknis
\item Diskusi dan evaluasi
\end{itemize} & 
UAS: Presentasi dan demo proyek akhir \\
\hline

\end{longtable}

\section{Pemetaan Pertemuan ke Bab dan Section}

\begin{table}[h]
\centering
\caption{Pemetaan RPS ke Buku Ajar}
\label{tab:pemetaan_rps_bab}
\small
\begin{tabular}{|c|l|l|}
\hline
\textbf{Pert.} & \textbf{Bab} & \textbf{Section} \\
\hline
1 & Bab 1 & 1.1-1.5 (Konsep Dasar Game Development) \\
2 & Bab 2 & 2.1-2.5 (Pengenalan Unity Engine) \\
3 & Bab 3 & 3.1-3.3 (Dasar Pemrograman C\#) \\
4 & Bab 4 & 4.1-4.5 (GameObject, Component, Prefab) \\
5 & Bab 5 & 5.1-5.6 (Input System dan Player Controller) \\
6 & Bab 6 & 6.1-6.6 (Physics, Collision, Rigidbody) \\
7 & Bab 7 & 7.1-7.5 (Animation dan Animator) \\
8 & UTS & Review Bab 1-7 \\
9 & Bab 8 & 8.1-8.5 (UI/UX Game) \\
10 & Bab 9 & 9.1-9.4 (Audio dalam Game) \\
11 & Bab 10 & 10.1-10.5 (Level Design dan Game Mechanics) \\
12 & Bab 11 & 11.1-11.5 (AI Sederhana untuk Game) \\
13 & Bab 12 & 12.1-12.5 (Optimization dan Debugging) \\
14 & Bab 13 & 13.1-13.3 (Build dan Deployment) \\
15 & Bab 13 & 13.4-13.5 (Final Project) \\
16 & UAS & Presentasi Proyek Akhir \\
\hline
\end{tabular}
\end{table}

\section{Metode Pembelajaran}

\begin{enumerate}
\item \textbf{Praktikum}: Hands-on practice menggunakan Unity Editor dan Visual Studio
\item \textbf{Project-Based Learning}: Pengembangan proyek game secara bertahap
\item \textbf{Problem-Based Learning}: Penyelesaian studi kasus dan permasalahan nyata
\item \textbf{Studi Kasus}: Analisis game-game populer dan mekaniknya
\item \textbf{Presentasi dan Diskusi}: Presentasi progress proyek secara berkala
\item \textbf{Peer Review}: Feedback antar mahasiswa untuk meningkatkan kualitas
\end{enumerate}

\section{Sistem Penilaian}

\begin{table}[h]
\centering
\caption{Distribusi Bobot Penilaian}
\label{tab:bobot_penilaian_rps}
\begin{tabular}{|l|c|}
\hline
\textbf{Komponen Penilaian} & \textbf{Bobot (\%)} \\
\hline
Tugas Praktikum (8 tugas) & 25\% \\
Kuis (2 kuis) & 10\% \\
UTS (Ujian Tengah Semester) & 20\% \\
Proyek Akhir & 30\% \\
UAS (Ujian Akhir Semester) & 15\% \\
\hline
\textbf{Total} & \textbf{100\%} \\
\hline
\end{tabular}
\end{table}

\section{Referensi}

Semua referensi menggunakan sumber open-access. Lihat Daftar Pustaka di akhir buku untuk detail lengkap.

\section{Indikator Kinerja Pembelajaran}

Setiap CPMK memiliki indikator kinerja yang terukur:
\begin{itemize}
\item \textbf{CPMK1}: Mampu menjelaskan minimal 5 genre game dan karakteristiknya
\item \textbf{CPMK2}: Mampu membuat skrip C\# dengan minimal 3 fungsi lifecycle
\item \textbf{CPMK3}: Mampu membuat scene dengan minimal 5 GameObject berbeda
\item \textbf{CPMK4}: Mampu implementasi player controller yang responsif
\item \textbf{CPMK5}: Mampu membuat simulasi fisika dengan minimal 3 objek interaktif
\item \textbf{CPMK6}: Mampu membuat animasi dengan minimal 3 states berbeda
\item \textbf{CPMK7}: Mampu membuat UI dengan minimal HUD dan menu
\item \textbf{CPMK8}: Mampu integrasi audio dengan minimal SFX dan BGM
\item \textbf{CPMK9}: Mampu merancang level dengan prinsip level design
\item \textbf{CPMK10}: Mampu implementasi AI dengan minimal 2 behavior patterns
\item \textbf{CPMK11}: Mampu optimasi game hingga mencapai target FPS
\item \textbf{CPMK12}: Mampu build dan deploy ke minimal 1 platform
\item \textbf{CPMK13}: Mampu membuat game utuh yang fungsional dan menarik
\end{itemize}
